%% ===================================================================
\section{Adversarial Equilibrium Theory}
\label{sec:theory}
%% ===================================================================

We now derive the theoretical machinery for the three questions posed in \cref{sec:background}. The main result (\cref{thm:phase_transition}) characterizes how vulnerability evolves with perturbation magnitude, and the constrained sensitivity matrix $S_c$ provides the unified object from which rankings, attacks, and radii follow.

\begin{theorem}[Vulnerability Characterization for Contractive Implicit GNNs]
\label{thm:phase_transition}
Let $F_\theta(Z, A) = \sigma(AZW^\top + X_{\mathrm{proj}})$ be an IGNN-class operator satisfying:
\begin{enumerate}
\item[(A1)] $\sigma$ is ReLU (or any 1-Lipschitz activation with $\sup|\sigma'| = 1$);
\item[(A2)] $W$ is spectral-norm constrained: $\norm{W}_2 \leq c$;
\item[(A3)] The state Jacobian satisfies $\norm{J_z}_2 \leq \kappa < 1$ at the fixed point $\zstar$ (contraction in operator norm).
\end{enumerate}
Define the critical budget
\begin{equation}
\ecrit = \frac{1 - \kappa}{\norm{W}_2}.
\end{equation}
Then structural perturbations $\delta A$ with $\norm{\delta A}_F = \varepsilon$ exhibit three regimes:

\noindent\textbf{(a) Subcritical} ($\varepsilon < \ecrit$): The perturbed operator remains contractive with a unique fixed point satisfying
\begin{equation}
\norm{\Delta \zstar}_F \leq \sigma_1(S) \cdot \varepsilon + O(\varepsilon^2),
\label{eq:shift_bound}
\end{equation}
where $S = (I - J_z)^{-1} J_A$ is the structural sensitivity matrix and $\sigma_1(S) \leq \norm{J_A}_{\mathrm{op}} / (1 - \kappa)$.

\noindent\textbf{(b) Critical} ($\varepsilon \to \ecrit$): The resolvent norm $\norm{(I - J_z')^{-1}}_2$ diverges as $\Omega(1/(\ecrit - \varepsilon))$. Intuitively, as the perturbation approaches the critical budget, the system's sensitivity to further perturbation grows without bound.

\noindent\textbf{(c) Supercritical} ($\varepsilon > \ecrit$): The sufficient contraction certificate becomes void ($\norm{J_z'}_2$ may exceed 1). The perturbed operator may still be contractive depending on perturbation direction, but the Banach fixed-point theorem no longer guarantees uniqueness, and the first-order guarantees from part~(a) cease to apply.
\end{theorem}

The intuition behind this three-regime picture is physical: a contractive system has a restoring force that pulls trajectories toward the equilibrium. Small perturbations shift the equilibrium smoothly (subcritical). As the perturbation approaches $\ecrit$, the restoring force weakens until the system becomes critically sensitive. Beyond $\ecrit$, the restoring force may vanish entirely, and the equilibrium can bifurcate or disappear.

\begin{proof}
\textbf{Part (a).} The Jacobian of $F$ with respect to $\mathrm{vec}(Z)$ takes the form $J_z = \mathrm{diag}(\sigma') \cdot (\Ahat \otimes W)$, where $\sigma' \in \{0, 1\}^{Nd}$ for ReLU (A1). The operator norm satisfies $\kappa = \norm{J_z}_2 \leq \norm{\Ahat}_2 \cdot \norm{W}_2 \cdot \sup|\sigma'| = \norm{\Ahat}_2 \cdot \norm{W}_2$. Applying the IFT to $G(z,A) = z - F(z,A) = 0$ at $(\zstar, A)$:
\begin{equation}
\Delta \zstar = (I - J_z)^{-1} J_A \cdot \mathrm{vec}(\delta A) + O(\norm{\delta A}^2).
\end{equation}
Taking operator norms and using the Neumann series bound $\norm{(I-J_z)^{-1}}_2 \leq 1/(1-\kappa)$ (convergent because $\kappa < 1$) yields~\eqref{eq:shift_bound}.

For contractivity preservation: the perturbed Jacobian satisfies $\norm{J_z'}_2 \leq (\norm{A}_2 + \norm{\delta A}_F) \cdot \norm{W}_2$ (using $\norm{\cdot}_2 \leq \norm{\cdot}_F$ for the perturbation). The condition $\varepsilon < (1-\kappa)/\norm{W}_2$ ensures $\norm{J_z'}_2 < 1$.

\textbf{Part (b).} As $\varepsilon \to \ecrit$, we have $\norm{J_z'}_2 \to 1$, so $\norm{(I - J_z')^{-1}}_2 \geq 1/(1 - \norm{J_z'}_2) \to \infty$. Unlike the spectral-radius bound $1/(1-\rho)$, which requires normality, the operator-norm bound holds for all $J_z$ with $\kappa < 1$. In our experiments, the pseudospectral index $\eta$ ranges from $1.02$ to $1.28$ (\cref{sec:convergence}), confirming mild non-normality.

\textbf{Part (c).} When $\norm{J_z'}_2 \geq 1$, the contraction mapping condition fails. The Banach theorem no longer guarantees existence or uniqueness; the iteration may converge to a different equilibrium, oscillate, or diverge.
\end{proof}

\textbf{Remark (conservativeness).} The critical budget $\ecrit$ is a sufficient, not necessary, condition for contractivity preservation. The perturbed system may remain contractive beyond $\ecrit$ if the perturbation is aligned with low-sensitivity directions. The unconstrained bound also uses worst-case $\norm{\cdot}_F$ relaxation. Under constrained perturbations, the first-order prediction achieves tightness $\approx 1.00$ (\cref{sec:cross_domain}). We use $\kappa = \norm{J_z}_2$ rather than the spectral radius $\rho(J_z)$ because the Neumann series requires operator-norm contractivity~\cite{stewart1990matrix}; for normal $J_z$ the two coincide, while for non-normal $J_z$ the $\kappa$-based bound is more conservative but requires no normality assumption.

\begin{proposition}[Optimal First-Order Structural Attack]
\label{prop:attack}
The perturbation $\delta A^*$ maximizing $\norm{\Delta \zstar}$ to first order subject to $\norm{\delta A}_F \leq \varepsilon$ is $\delta A^* = \varepsilon \cdot \mathrm{reshape}(v_1, N \times N)$, where $v_1$ is the leading right singular vector of $S$. The maximum first-order shift is $\varepsilon \cdot \sigma_1(S)$.
\end{proposition}

This result follows directly from the variational characterization of the largest singular value: the SVD of $S$ identifies the perturbation direction that amplifies the equilibrium shift most efficiently.

\textbf{Constrained sensitivity.} Under the threat model of \cref{sec:background} (symmetric, edge-only, $\norm{\delta \Ahat}_F \leq \varepsilon$), we construct the \emph{constrained sensitivity matrix} $S_c \in \R^{Nd \times |E|}$ with column
\begin{equation}
[S_c]_{:,k} = S_{:,iN+j} + S_{:,jN+i}
\end{equation}
for each edge $(i,j) \in E$. This reduces the perturbation space from $N^2$ dimensions to $|E|$ dimensions and enforces symmetry by construction. The constrained-optimal attack uses $\sigma_1(S_c)$, achieving tightness $\approx 1.00$ empirically (\cref{sec:cross_domain}). The construction of $S_c$ is the central technical contribution of this work: it transforms unconstrained sensitivity analysis into a tight, practical tool for realistic graph perturbations.

\begin{proposition}[Per-Node First-Order Sensitivity Radius]
\label{prop:radius}
Assume a linear classification head $f(z) = Wz + b$ (as in standard IGNN). For node $v$ with margin $m_v = f_{y_v}(\zstar_v) - \max_{c \neq y_v} f_c(\zstar_v) > 0$, the first-order sensitivity radius is
\begin{equation}
r_v = \frac{m_v}{\norm{W_{y_v} - W_{c^*}}_2 \cdot \norm{S_v}_2},
\label{eq:radius}
\end{equation}
where $S_v$ denotes the block-rows of $S$ corresponding to node $v$, $W_{y_v}$ and $W_{c^*}$ are the classification weight vectors for the predicted class $y_v$ and runner-up class $c^*$, respectively. Any perturbation with $\norm{\delta \Ahat}_F < r_v$ preserves the classification of node $v$ to first order.
\end{proposition}

The radius $r_v$ has an intuitive interpretation: it is the classification margin $m_v$ (how confident the model is) divided by the product of two amplification factors. The term $\norm{W_{y_v} - W_{c^*}}_2$ measures how sensitive the decision boundary is to hidden-state changes, while $\norm{S_v}_2$ measures how sensitive node $v$'s equilibrium position is to structural perturbation. Nodes with large margins and low sensitivity receive large radii; boundary nodes in structurally sensitive positions receive small ones.

\textit{Proof sketch.} By Theorem~1(a), $\Delta \zstar_v \approx S_v \cdot \mathrm{vec}(\delta A)$, so $|\Delta f_{y_v}| \leq \norm{\partial f / \partial \zstar_v} \cdot \norm{S_v} \cdot \norm{\delta A}_F$. Misclassification requires $|\Delta f_{y_v}| \geq m_v$; inverting gives $r_v$. Note $S_v$ already contains the $(I-J_z)^{-1}$ factor.

These radii are deterministic (not probabilistic like randomized smoothing~\cite{bojchevski2020efficient}) and first-order: locally tight for small $\varepsilon$ but increasingly conservative as perturbation magnitude grows. \cref{sec:adaptive} verifies empirically that they hold (0\% breach rate at $\varepsilon = 0.01$, $< 1\%$ at $\varepsilon = 0.10$).

\subsection{Generalization to Explicit GNNs}

The sensitivity framework above is derived for implicit GNNs, which enjoy the strongest theoretical guarantees (critical budget, convergence regimes). The $S_c$ construction itself, however, generalizes to any differentiable $K$-layer GNN as a computational tool for vulnerability analysis.

\begin{proposition}[Structural Sensitivity for $K$-Layer GNNs]
\label{prop:explicit}
Let $Z_K = g_K \circ \cdots \circ g_1(X, A)$ be a $K$-layer GNN where each layer $g_l$ is differentiable with respect to both its input and $A$. Define the \emph{unrolled sensitivity matrix}
\begin{equation}
S_K = \frac{\partial \mathrm{vec}(Z_K)}{\partial \mathrm{vec}(A)} = \sum_{l=1}^{K} \left(\prod_{k=l+1}^{K} J_z^{(k)}\right) J_A^{(l)},
\label{eq:unrolled}
\end{equation}
where $J_z^{(k)} = \partial g_k / \partial Z_{k-1}$ and $J_A^{(l)} = \partial g_l / \partial \mathrm{vec}(A)$. Then:

\noindent\textbf{(a)} The first-order shift bound $\norm{\Delta Z_K}_F \leq \sigma_1(S_K) \cdot \norm{\delta A}_F + O(\norm{\delta A}^2)$ holds, with
\begin{equation}
\sigma_1(S_K) \leq \sum_{l=1}^{K} \left(\prod_{k=l+1}^{K} \norm{J_z^{(k)}}_2\right) \norm{J_A^{(l)}}_2.
\end{equation}

\noindent\textbf{(b)} The constrained construction $S_c$ and per-node radius $r_v$ (\cref{prop:radius}) apply identically with $S_K$ in place of $S$.
\end{proposition}

For weight-tied models ($J_z^{(k)} = J_z$, $J_A^{(l)} = J_A$ for all $k, l$), the bound simplifies to $\sigma_1(S_K) \leq \norm{J_A}_2 \cdot (1 - \kappa^K)/(1-\kappa)$, which converges to the IGNN bound $\norm{J_A}_2 / (1-\kappa)$ as $K \to \infty$.

\textbf{What explicit GNNs gain and lack.} From \cref{prop:explicit}, any differentiable GNN inherits: (i) the $S_c$ construction, (ii) SVD-optimal attacks, (iii) per-edge vulnerability rankings, and (iv) per-node first-order sensitivity radii. What explicit models \emph{lack} is the critical budget $\ecrit$ and the three-regime convergence characterization of \cref{thm:phase_transition}, which require the contraction assumption (A3). \cref{sec:explicit_extension} validates this generalization on six explicit architectures.
